\ulohahead{specializace UI-SU}{Bayesovské sítě a pravděpodobnostní inference}
\begin{zadani}
\begin{enumerate}[label=(\alph*)]
  \item \bod{1} Definujte Bayesovskou síť a její vztah k úplné sdružené distribuci (faktorizace).
  \item \bod{2} Jak se síť konstruuje (podmíněná nezávislost) a jak se počítá exaktní inference enumerací?
  \item \bod{3} Popište eliminaci proměnných a aproximační inferenci (Monte Carlo: zamítání, vážení věrohodností).
\end{enumerate}
\end{zadani}

\begin{reseni}
\textbf{(a)} \emph{Bayesovská síť} = orientovaný acyklický graf, kde uzly jsou náhodné proměnné a hrany vyjadřují přímou závislost; každý uzel má tabulku podmíněných pravděpodobností $P(X_i\mid \mathrm{rodiče}(X_i))$. Kóduje \emph{úplnou sdruženou distribuci} faktorizací
\[
P(X_1,\dots,X_n)=\prod_{i=1}^n P\big(X_i\mid \mathrm{rodiče}(X_i)\big),
\]
což šetří parametry oproti plné tabulce (využívá podmíněné nezávislosti).

\textbf{(b)} \emph{Konstrukce}: seřaď proměnné (ideálně příčiny před následky) a pro každou $X_i$ vyber za rodiče $\mathrm{Pa}(X_i)$ takovou podmnožinu předchozích proměnných, aby byla $X_i$ podmíněně nezávislá na ostatních předchozích proměnných při daných $\mathrm{Pa}(X_i)$. \emph{Inference enumerací} dotazu $P(Q\mid e)$: $P(Q\mid e)\propto\sum_{\text{skryté } h}P(Q,e,h)$, kde sčítáme součiny faktoru z faktorizace přes všechny hodnoty skrytých proměnných a nakonec normujeme.

\textbf{(c)} \emph{Eliminace proměnných} sčítá zprava efektivněji: postupně ,,vysčítá'' skryté proměnné, mezivýsledky drží jako faktory a opakované součiny nepřepočítává - výrazně rychlejší než holá enumerace. \emph{Aproximace (Monte Carlo)}: \emph{zamítavé vzorkování} (rejection) generuje vzorky z priorů a zahodí ty neslučitelné s $e$ (neefektivní při vzácném $e$); \emph{vážení věrohodností} (likelihood weighting) fixuje pozorované proměnné a každý vzorek váží součinem jejich pravděpodobností - žádný vzorek nezahazuje.
\end{reseni}

\begin{jadro}
Bayesovská síť = DAG + tabulky $P(X_i\mid\text{rodiče}(X_i))$, kóduje $P(X_1,\dots,X_n)=\prod_i P(X_i\mid\text{rodiče}(X_i))$. Inference: $P(Q\mid e)\propto\sum_{\text{skryté }h}P(Q,e,h)$ (enumerace, efektivněji eliminace proměnných).
\end{jadro}

\kalibrace{Bayesovské sítě (~22\,\% s Bayesovským učením) jsou opakovaná ,,Základy UI'' otázka. Faktorizace + enumerace = 2-bodová záloha; tady je AI rodina, kde je vhodné mít i hlubší zvládnutí.}
\past{Faktorizace je součin $P(X_i\mid\text{rodiče})$ - ne marginálu. Likelihood weighting nezahazuje vzorky (na rozdíl od rejection), proto je efektivnější při vzácné evidenci.}
